\chapter{Methodology}
\label{ch:methodology}

\section{The analysis in six steps}

The analysis has six ordered steps. A measure reaches the economic tests only
after the fixed statistical comparison, preventing a later return result from
replacing the measure chosen at the start.

\begin{enumerate}[leftmargin=1.0cm]
  \item Compare nine firm-day summaries using one outcome and one multiplicity
  correction.
  \item Test whether the selected measure adds to average sentiment.
  \item Split the return window into parts to see where the relationship appears.
  \item Apply the unchanged baseline once to the fixed 2020--2023 testing period.
  \item Translate the coefficient into return scale, portfolio performance,
  turnover and costs.
  \item Evaluate a same-session FNSPID risk diagnostic, then repeat the
  coefficient and economic work in the LSEG data.
\end{enumerate}

No statistic covers association, timing and practical value. Rank correlations
measure co-movement. Portfolio returns require turnover and costs. A drawdown
claim also requires matched market exposure and verified signal availability.

The selected coefficient is exploratory because the aggregation analysis
changed before the final research question was set. Comparing its $p$-value with
$0.05/9=0.00556$ provides a conservative nine-test check, but does not capture
all uncertainty from choosing the best rule \citep{harvey2016multiple}. Some
later checks were also designed after seeing the association.

The testing-period estimates, test groups and decision rule were fixed before they
were calculated. However, some outcomes from those years had already been
examined.

\section{Step 1: nine firm-day aggregation rules}

Let $\mathcal{J}_{it}$ contain the stories assigned to firm $i$ on session $t$,
with $n_{it}=|\mathcal{J}_{it}|$. The continuous score $s_{ijt}$ comes from
\cref{eq:continuous-score}, and $h_{ijt}\in\{-1,0,1\}$ is the hard negative,
neutral or positive label. Eight rules use only the current firm-day. The ninth
carries a decayed state through time.

\begin{enumerate}[leftmargin=0.9cm]
  \item \textbf{Mean hard label:}
  $a^{\mathrm{hard}}_{it}=n_{it}^{-1}\sum_j h_{ijt}$.

  \item \textbf{Mean continuous score:}
  $a^{\mathrm{mean}}_{it}=n_{it}^{-1}\sum_j s_{ijt}$.

  \item \textbf{Median continuous score:}
  $a^{\mathrm{median}}_{it}=\operatorname{median}_{j}(s_{ijt})$.

  \item \textbf{Ten-per-cent trimmed mean:} sort the scores, drop
  $k=\lfloor0.1n_{it}\rfloor$ observations from each tail and average the rest.
  With small $n_{it}$, $k=0$ and the rule collapses to the mean.

  \item \textbf{Negative-story share:}
  \begin{equation}
  q^{-}_{it}=\frac{1}{n_{it}}\sum_{j\in\mathcal{J}_{it}}
  \mathbf{1}(h_{ijt}=-1).
  \label{eq:negative-share}
  \end{equation}

  \item \textbf{Dispersion:} the sample standard deviation of $s_{ijt}$, set to
  zero on a one-story firm-day.

  \item \textbf{Strongest event:} prefer a non-recap story, then take the
  remaining story with the largest $|s_{ijt}|$. An exact tie in opposite
  directions is recorded as missing, and the return never breaks a tie.

  \item \textbf{Log-count attention:}
  $a^{\mathrm{attn}}_{it}=a^{\mathrm{mean}}_{it}\log(1+n_{it})$. This rule picks
  out no particular story. It gives more cross-sectional weight to a mean built
  from more stories.

  \item \textbf{Decayed state:} convert the current continuous mean into an
  impulse $u_{it}=\operatorname{sign}(a^{\mathrm{mean}}_{it})|a^{\mathrm{mean}}_{it}|^{1.5}$.
  Writing $d_{it^{-}}$ for the state at firm $i$'s previous news-bearing session
  and $\Delta_{it}$ for the number of exchange sessions since it,
  \begin{equation}
   d_{it}=2^{-\Delta_{it}/3}d_{it^{-}}
   \Bigl(1-\tfrac{3}{4}\mathbf{1}\bigl[u_{it}d_{it^{-}}<0\bigr]\Bigr)+u_{it}.
   \label{eq:decayed-state}
  \end{equation}
  Tone persists across sessions, and a strictly opposite non-zero impulse
  resets three quarters of the decayed prior state before the new impulse enters.
  The state is initialised at zero for each firm's first news-bearing session.
\end{enumerate}

The rules are deliberately simple: learning an aggregator would add another
selection problem and obscure what the mean discarded. The candidates cover a
threshold, dispersion, trimmed extremes, story selection, attention and a
cross-session state.

For each session, firms are ranked by the signal and by return. The Spearman
correlation asks whether firms with higher signal ranks also tend to have higher
or lower return ranks. For rule $a$, it is

\begin{equation}
 IC^a_t=\operatorname{corr}_{S}\left(a_{it},r^A_{i,t+1}\right).
 \label{eq:daily-ic}
\end{equation}

The reported coefficient is $\overline{IC}^a=T^{-1}\sum_t IC^a_t$, requiring at
least three complete firms per session. Five-lag Newey-West HAC errors allow
changing variance and short-run dependence \citep{newey1987}. The nine two-sided
$p$-values form one Benjamini-Hochberg family at $q=0.05$
\citep{benjamini1995}. Story-count groups are descriptive. Expected directions
are fixed: positive-score summaries point upward, while negative share and
dispersion point downward.

\section{Step 2: the relationship after controlling for the mean}

For each session, firms are ranked by return, negative share, average sentiment
and story count. The regression asks whether negative-share rank is still
related to return rank after accounting for the other two news measures. The
formal calculation uses within-session percentile ranks, with tied values given
their average rank and each rank shifted by one half. Let $N_t$ be the available
firms, $L$ the number strictly below a firm on $x$, and $E$ its tie-group size:

\begin{equation}
 \widetilde x_{it}=R_t(x_{it})-\tfrac12,
 \qquad
 R_t(x_{it})=\frac{L+\tfrac12(E+1)}{N_t}.
 \label{eq:shifted-rank}
\end{equation}

\noindent Without ties, shifted ranks span $N_t^{-1}-\tfrac12$ to $\tfrac12$.
Their mean is $1/(2N_t)$, which the daily intercept absorbs. Because ties narrow
the support, effect translations use observed rank spreads. The base model is

\begin{equation}
 \widetilde r^A_{i,t+1}=\alpha_t+\beta_t\widetilde q^{-}_{it}
 +\gamma_t\widetilde{\bar s}_{it}
 +\delta_t\widetilde{\log(1+n_{it})}+\varepsilon_{it},
 \label{eq:conditional-base}
\end{equation}

\noindent where $\bar s_{it}$ is average continuous sentiment and tildes denote
shifted ranks. The reported coefficient averages the daily slopes:

\begin{equation}
 \bar\beta=\frac{1}{T}\sum_{t=1}^{T}\beta_t.
 \label{eq:mean-beta}
\end{equation}

A session needs ten complete firms and an estimable design matrix. The slope
series uses a five-lag HAC error. Ranking limits outliers and gives FNSPID and
LSEG a common scale. Since ties prevent a full unit change in negative-share
rank, translations use the observed spread.

The price-path specification adds three shifted ranks to test the reversal
finding of \citet{chan2003news}:

\begin{equation}
 \widetilde r^A_{i,t+1}=\alpha_t+\beta_t\widetilde q^{-}_{it}
 +\gamma_t\widetilde{\bar s}_{it}+\delta_t\widetilde{\log(1+n_{it})}
 +\theta_{1t}\widetilde r^{(1)}_{it}+\theta_{5t}\widetilde r^{(5)}_{it}
 +\phi_t\widetilde\sigma^{(20)}_{it}+\varepsilon_{it}.
 \label{eq:conditional-full}
\end{equation}

\noindent Here $r^{(1)}$ is the adjusted-open return from $t-1$ to $t$,
$r^{(5)}$ is the cumulative return from $t-5$ to $t$, and $\sigma^{(20)}$ is the
standard deviation of the twenty returns ending at formation. Every control
ends before the outcome starts. The full model is the main reversal test.
Intermediate rows show the effect of adding controls.

Three post-result measurement checks form a separate Benjamini--Hochberg family.
The first uses a close-to-close market-adjusted outcome. The second further
removes the component explained by a trailing Fama--French three-factor (FF3)
model \citep{famafrench1993}. The third keeps the primary outcome but replaces
hard negative share with average $p(\mathrm{negative})$. Factor betas use the
252 sessions ending at formation, require 126 observations and exclude the
outcome session.

The rank control is linear. On one-story days, negative share is a thresholded
class from the same probability vector as mean sentiment, so its coefficient may
capture a nonlinear boundary rather than aggregation. The result is conditional
only on the stated controls. A more flexible model would be another
training-sample test and is left for a prospectively specified study.

On price-complete rows, the pooled 90th-minus-10th percentile negative-share-rank
spread is 0.56056. The mean session spread is 0.52722. Multiplying either by
$100\bar\beta$ gives return-rank percentile points. A hypothetical change of
0.8 exceeds the observed pooled range and is not used.

\section{Step 3: locating the association inside the return window}

The timing test splits stock and SPY returns into assigned-open-to-close,
assigned-close-to-next-open and previous-close-to-assigned-open parts, plus a
separate previous-open-to-assigned-open window. Gross-return identities are
checked to $10^{-12}$, and the three new training legs form one correction
family. The testing-period results and contrasts between periods are additional
checks rather than pre-specified tests. Without story arrival times, the split
cannot establish sequence.

A second check uses a 252-session rolling stock-SPY beta ending at the assigned
open and requiring 126 observations. One model beta-adjusts the outcome. Another
also controls for beta rank. The two tests are corrected within each period. A
separate correction is applied to the two between-period contrasts. Neither
analysis estimates causality or tradable alpha.

\section{Step 4: the fixed 2020--2023 test}

The fixed testing procedure evaluates two pre-specified families once. The first
repeats the nine two-sided information-coefficient tests from
\cref{eq:daily-ic}. The second repeats \cref{eq:conditional-base} on all firm-days
and the pre-specified multi-story subset, with a shared correction and negative
expected direction. It does not examine testing-period portfolios, alternative
outcomes, price controls, thresholds, event types, scorers or other story-count
groups.

Before the two fixed 2020--2023 coefficient families were calculated, an
amendment added prospective power information and a direct
training-versus-testing contrast.
Expected precision scales the training HAC error by $\sqrt{2,264/998}$ and uses
the corrected two-test threshold, a two-sided size of 0.025. A descriptive check
uses the observed testing error instead.

To compare the periods, I stack the unchanged daily coefficient series and estimate

\begin{equation}
 \beta_t=\alpha+\Delta D^{\mathrm{test}}_t+u_t
 \label{eq:temporal-contrast}
\end{equation}

\noindent with HAC(5) uncertainty. $D^{\mathrm{test}}_t$ is zero in
training and one in testing, so $\alpha=\bar\beta_{\mathrm{training}}$ and
$\Delta=\bar\beta_{\mathrm{testing}}-\bar\beta_{\mathrm{training}}$. The contrast
uses \cref{eq:conditional-base}, excluding the price-path model. It is needed
because significance in one period and non-significance in another do not prove
a difference. It does not change the verdict from the primary tests.

\section{Step 5: translation into economic terms}

This step covers the double-sort comparison, portfolio costs and learned
no-trade filters.

\subsection{Double-sort comparison}

The double sort first groups firms with similar average sentiment. Within each
group, it then compares firms with low, middle and high negative-story shares.
The high-minus-low difference shows whether negative share matters among firms
with similar average sentiment. Fractional ties prevent symbol order or returns
from breaking them. Cells need five firms and sessions need three usable groups.
The daily high-minus-low series uses HAC(5) uncertainty and is repeated for
multi-story days. Its $p$-value is descriptive. Earlier fixed cut-offs produced
no usable sessions, so the tie-aware replacement was fixed before estimation.

\subsection{Rank portfolios, turnover and costs}

Each rule also forms a daily cross-sectional rank portfolio. Expected-direction
ranks are scaled to equal long and short gross exposure. Positions formed at
open $t$ earn abnormal returns to open $t+1$. Let $w_{it}$ be the target weight,
$r_{it}$ the raw stock return into open $t$, and $w^{\mathrm{pre}}_{it}$ the
previous target after that return:

\begin{equation}
 w^{\mathrm{pre}}_{it}=
 \frac{w_{i,t-1}(1+r_{it})}
 {1+\sum_j w_{j,t-1}r_{jt}}.
 \label{eq:weight-drift}
\end{equation}

Turnover is reported one-way, as half the traded notional:

\begin{equation}
 \mathrm{TO}_t=\tfrac{1}{2}\sum_i\left|w_{it}-w^{\mathrm{pre}}_{it}\right|.
 \label{eq:turnover}
\end{equation}

A quoted per-side cost applies to both sales and purchases, giving a charge of
$2c\,\mathrm{TO}_t$ and a net return of

\begin{equation}
 r^{\mathrm{net}}_{p,t+1}=\sum_i w_{it}r^A_{i,t+1}-2c\,\mathrm{TO}_t.
 \label{eq:portfolio-net}
\end{equation}

Starting a dollar-neutral portfolio with gross exposure of one gives
$\mathrm{TO}_t=0.5$ and costs $c$. Final liquidation charges the same rate on
half the remaining absolute notional. The main cost is 10 basis points per side.
Descriptive break-even cost is
$\bar r^{\mathrm{gross}}_p/(2\overline{\mathrm{TO}})$, annualised turnover is
$252\overline{\mathrm{TO}}$, and Sharpe ratios use the same annualisation. A
non-positive denominator in \cref{eq:weight-drift} implies insolvency and leaves
later cost metrics undefined unless returns are capped.

\subsection{Learned no-trade filters}

Learned no-trade thresholds test sensitivity to the simple rule. Logistic
regression, gradient boosting and a multilayer perceptron (MLP) predict whether
to trade. The comparison covers testing-period discrimination, measured by the
area under the receiver operating characteristic curve (AUC), trading activity
and net return against a fixed band, an activity-matched rule and cash. A network
fitted to shuffled labels is the negative control.

\section{Step 6: same-session risk tests}

The risk analysis asks whether weak directional information coincides with
periods when lower exposure would have helped. The FNSPID exercise is a
same-session diagnostic, not a strategy test based on information known in
advance, because the data do not show when each story became available.

\subsection{Volatility benchmark}

The benchmark reduces market exposure when recent volatility points to higher
risk. The news rule reduces exposure further on unusually negative-news days.
It adds value only if it improves downside outcomes beyond the volatility rule
at a similar level of market exposure.

The sentiment-free HAR benchmark forecasts the average squared SPY return over
the next five sessions. At decision open $t$,
$u_t=r^2_{\mathrm{SPY},t-1\rightarrow t}$ is the previous session's squared
adjusted-open return and is already known. The forecast target is
$v_t^{(5)}=5^{-1}\sum_{j=0}^{4}r^2_{\mathrm{SPY},t+j\rightarrow t+j+1}$. With
$\epsilon=10^{-8}$, training-period ordinary least squares (OLS) estimates

\begin{equation}
 \log(v_t^{(5)}+\epsilon)=\alpha+\beta_1\log(u_t+\epsilon)
 +\beta_5\log\left(\frac{1}{5}\sum_{j=0}^{4}u_{t-j}+\epsilon\right)
 +\beta_{22}\log\left(\frac{1}{22}\sum_{j=0}^{21}u_{t-j}+\epsilon\right)+e_t.
 \label{eq:har-variance}
\end{equation}

Let $x_t$ collect the three logged volatility measures in
\cref{eq:har-variance}, excluding the intercept. Converting the forecast back
from logs gives
$\widehat v_t=\kappa\exp(\widehat\alpha+x_t'\widehat\beta)$, where the positive
training-period scale correction is
$\kappa=\operatorname{mean}\{v_t^{(5)}/\exp(\widehat\alpha+x_t'\widehat\beta)\}$;
the mean is taken over training-period targets. This correction adjusts the
average forecast level after the conversion from logs.

The fixed test model is fitted only to five-session targets whose windows end by
31 December 2019. A historical sensitivity starts after 504 rows and refits
expanding OLS every 21 sessions using completed target windows. Forecasts are
converted into exposure as
$\operatorname{clip}\{0.10/\sqrt{252\widehat v_t},0.25,1\}$, so higher forecast
volatility reduces market exposure. Exposure cannot fall below 0.25 or exceed
1. Cash earns zero, costs are 2 basis points per side, and the overlay cannot
exceed HAR exposure.

\subsection{Same-session news rule}

The sentiment input is market-wide negative-news pressure: the
headline-count-weighted negative share, compared with its recent history.

\begin{equation}
 P_t=\frac{\sum_i q^{-}_{it}n_{it}}{\sum_i n_{it}},
 \qquad
 z_t=\frac{P_t-\mu_{t-1}}{\sigma_{t-1}},
 \label{eq:negative-pressure}
\end{equation}

\noindent Here $\mu_{t-1}$ and $\sigma_{t-1}$ use the 252 sessions ending at
$t-1$ and require 126 observations. Current-session news enters $P_t$ and the
numerator of $z_t$. Only the normalisation is historical. The signal is therefore
not lagged. A two-threshold state defines the exposure modifier:

\begin{equation}
 S_t=\begin{cases}
 1 & \text{if } S_{t-1}=0 \text{ and } z_t\geq 1.5,\\
 0 & \text{if } S_{t-1}=1 \text{ and } z_t\leq 0.5,\\
 S_{t-1} & \text{otherwise,}
 \end{cases}
 \label{eq:hysteresis}
\end{equation}

\noindent with $S_0=0$ and the prior state retained when $z_t$ is unavailable.
When $S_t=1$, HAR exposure is multiplied by 0.25. Separate entry and exit
thresholds limit switching. The state calculated for session $t$ sets exposure
in that same session; it is not delayed by one session. The result therefore
cannot be interpreted as a strategy based on information known before the
session began.

\subsection{Matched controls and uncertainty}

For a like-for-like comparison, a constant multiplier rescales HAR to the news
rule's mean 2020--2023 exposure. This scale is calculated from the testing
results, so it provides an after-the-fact explanation of lower exposure rather
than a control selected in training.

Paired block bootstraps keep nearby sessions in fixed twenty-session blocks and
wrap them at the sample boundary, using 4,999 or 9,999 repetitions as recorded
in each manifest \citep{kunsch1989}. Return differences and reductions in
squared negative returns are separate outcomes. Circular shifts test unusual
alignment of the same-session risk schedule with downside outcomes, but cannot
establish signal availability. Additional checks omit each episode and
February--April 2020 in turn.

\section{The LSEG analysis and the prompt experiments}

The unchanged FNSPID pressure rule is first applied to both LSEG periods. A
restricted search then tests other aggregate rules under a fixed plan. The
selected aggregate rule is compared with HAR and a constant multiplier chosen
in training.

A separate firm-level rule is compared with full investment and with fixed
company weights that match its training exposure. A rule passes only if it
improves tail risk or utility, retains an acceptable lower bound on return and
is stable across checks. Half-sample and leave-one-firm-out results show whether
the finding depends on a small part of the data.

Prompt experiments compare generic tone with firm-specific economic labels.
Five Gemma 4 26B prompts target one-session reaction, delayed reaction,
cash-flow revision, surprise catalysts and a red-team consensus view. A sixth
targets the first tradable price reaction. Each returns a score on a fixed grid.
Reuters headlines and company names passed through OpenRouter to DeepInfra under
recorded no-retention controls. Returns were excluded from every payload.

Training uses 2024. A prompt is retained only if it meets five pre-specified
conditions. It must have a positive net mean after 10 basis points per
side, a positive net Sharpe and a break-even cost of at least 10 basis points.
It must also be positive in three quarters and beat the same-date Gemma baseline
in three quarters. If none meets all five conditions, the procedure computes no
2025 prompt portfolio.
\ifdefined\withoutappendices
\else
  \Cref{app:prompts} records the tested wording because the null depends on what
  the model was asked.
\fi

\section{Reporting null results}

A null result may mean that there is little effect, or that the sample is too
small to detect it. The minimum detectable effect (MDE) is the smallest effect
that the test has an 80\% chance of detecting under the stated assumptions. For
standard error $SE$, the nominal two-sided 5\% MDE is

\begin{equation}
 MDE_{0.80}=\left(z_{0.975}+z_{0.80}\right)SE=2.8016SE.
 \label{eq:mde}
\end{equation}

For $m$ tests, the conservative first-BH threshold replaces $z_{0.975}$ with
$z_{1-0.05/(2m)}$. This is sufficient for the first ordered test, not exact
adaptive-BH power. MDEs are reported for LSEG transfer, story-type,
earnings-window, publisher and learned-threshold nulls. The analysis records
families specified before the main analysis and labels later families separately.
